\documentclass{article}
\usepackage[utf8]{inputenc}
\usepackage{amsmath,amssymb}
\usepackage[most]{tcolorbox}
\usepackage{geometry}
\geometry{margin=2.5cm}

\newcommand{\step}[1]{\par\noindent\hangindent=3.5em\hangafter=1%
  \makebox[3.5em][l]{\textbf{#1.}}}
\newcommand{\steptag}[1]{\hfill{\small\textsf{[#1]}}}

\newtcolorbox{proofbox}[1]{
  colback=gray!5, colframe=gray!60,
  fonttitle=\bfseries, title={#1},
  breakable, enhanced,
  left=4pt, right=4pt, top=4pt, bottom=4pt,
  before skip=10pt, after skip=10pt
}

\begin{document}

\begin{proofbox}{The signed-idempotent and stochastic-idempotent formulati...}
\step{1} The signed-idempotent and stochastic-idempotent formulations of classical stability are equivalent up to universal constants: Q row-stochastic with ||Q\textasciicircum{}2-Q|| <= eta gives P=theta(2Q-1) signed affine retraction with ||P-Q|| <= C eta and neg mass delta <= C eta, and conversely row-normalising p\_i\textasciicircum{}+ gives Q with ||P-Q|| <= 2 delta, ||Q\textasciicircum{}2-Q|| <= 6 delta+4 delta\textasciicircum{}2. \steptag{validated} \\[6pt]
\step{1.1} Forward direction: for sufficiently small eta, any row-stochastic Q with ||Q\textasciicircum{}2-Q|$|\_\{inf-\rangle$inf\} <= eta produces an exact signed idempotent P=theta(2Q-I) with P1=1, P\textasciicircum{}2=P, ||P-Q|$|\_\{inf-\rangle$inf\} <= C eta, and negative mass delta(P) <= C eta for a universal C. \steptag{validated} \\[6pt]
\step{1.1.1} Let X=2Q-I. Since Q is row-stochastic, X1=1, ||X|$|\_\{inf-\rangle$inf\} <= 3, and X\textasciicircum{}2-I=4(Q\textasciicircum{}2-Q), so ||X\textasciicircum{}2-I|$|\_\{inf-\rangle$inf\} <= 4 eta. \steptag{validated} \\[4pt]
\step{1.1.2} Using the validated sibling 1.1.1, write Y=X\textasciicircum{}2-I with X1=1, ||X|$|\_\{inf-\rangle$inf\}<=3, and ||Y|$|\_\{inf-\rangle$inf\}<=4 eta. For eta<=1/8 define B=(X\textasciicircum{}2)\textasciicircum{}(-1/2)=(I+Y)\textasciicircum{}(-1/2) by the convergent binomial series and set S=XB and theta(X)=(I+S)/2. Then B commutes with X, B1=1, ||B-I|$|\_\{inf-\rangle$inf\}<=8 eta; hence S\textasciicircum{}2=I, S1=1, and ||S-X|$|\_\{inf-\rangle$inf\}<=24 eta. \steptag{validated} \\[6pt]
\step{1.1.2.1} From 1.1.1, Y=X\textasciicircum{}2-I satisfies Y1=X\textasciicircum{}2 1-1=X1-1=0, ||Y|$|\_\{inf-\rangle$inf\}<=4 eta, and ||X|$|\_\{inf-\rangle$inf\}<=3. If eta<=1/8 then ||Y||<=1/2<1, so the binomial series B=sum\_\{m>=0\} binom(-1/2,m)Y\textasciicircum{}m converges absolutely in operator norm. Because B is the norm limit of polynomials in Y=X\textasciicircum{}2-I, it commutes with X. Since Y1=0, Y\textasciicircum{}m1=0 for all m>=1, hence B1=1. \steptag{validated} \\[4pt]
\step{1.1.2.2} For t=||Y|$|\_\{inf-\rangle$inf\}<=1/2, the absolute coefficient sum gives ||B-I||<=sum\_\{m>=1\}|binom(-1/2,m)|t\textasciicircum{}m=(1-t)\textasciicircum{}(-1/2)-1<=2t. With t<=4 eta this yields ||B-I|$|\_\{inf-\rangle$inf\}<=8 eta. \steptag{validated} \\[4pt]
\step{1.1.2.3} The scalar identity ((1+z)\textasciicircum{}(-1/2))\textasciicircum{}2(1+z)=1 on |z|<1 transfers to the absolutely convergent power series in Y; therefore B\textasciicircum{}2(I+Y)=I, i.e. B\textasciicircum{}2 X\textasciicircum{}2=I. Since B commutes with X, S=XB satisfies S\textasciicircum{}2=XBXB=X\textasciicircum{}2B\textasciicircum{}2=I. \steptag{validated} \\[4pt]
\step{1.1.2.4} Using B1=1 and X1=1 from 1.1.2.1, S1=XB1=X1=1. Also S-X=X(B-I), so ||S-X|$|\_\{inf-\rangle$inf\}<=||X|$|\_\{inf-\rangle$inf\}||B-I|$|\_\{inf-\rangle$inf\}<=3*8 eta=24 eta by 1.1.1 and 1.1.2.2. Thus the amended statement holds with K=8 and C=24. \steptag{validated} \\[6pt]
\step{1.1.2.4.1} Explicit dependency import for node 1.1.2.4: this algebraic step depends on the validated local facts X1=1 and B1=1 from 1.1.2.1, ||X|$|\_\{inf-\rangle$inf\}<=3 from 1.1.1, and ||B-I|$|\_\{inf-\rangle$inf\}<=8 eta from 1.1.2.2. These are recorded as required validated dependencies here, so no unvalidated sibling fact is assumed inside this node. \steptag{archived} \\[4pt]
\step{1.1.2.4.2} Conditional algebra/QED: with S=XB and the imported identities B1=1 and X1=1, S1=XB1=X1=1. Also S-X=XB-X=X(B-I); by submultiplicativity of ||.|$|\_\{inf-\rangle$inf\} and the imported bounds, ||S-X|$|\_\{inf-\rangle$inf\}<=||X|$|\_\{inf-\rangle$inf\}||B-I|$|\_\{inf-\rangle$inf\}<=3*8 eta=24 eta. Thus node 1.1.2.4 proves exactly the consequences it claims, with the contested inputs supplied only through the explicit validated dependencies. \steptag{archived} \\[4pt]
\step{1.1.2.4.3} Dependency-recorded QED: relying explicitly on validated nodes 1.1.1, 1.1.2.1, and 1.1.2.2, we may use X1=1, B1=1, ||X|$|\_\{inf-\rangle$inf\}<=3, and ||B-I|$|\_\{inf-\rangle$inf\}<=8 eta. Then S1=XB1=X1=1 and S-X=X(B-I), so ||S-X|$|\_\{inf-\rangle$inf\}<=3*8 eta=24 eta by submultiplicativity. This child supplies the missing dependency relationship challenged in ch-c6006369fb6ad4bf. \steptag{validated} \\[4pt]
\step{1.1.2.5} Dependency bridge/QED: node 1.1.2 is not derived from the bare inequality ||X\textasciicircum{}2-I||<1. Its imported hypotheses X1=1, ||X|$|\_\{inf-\rangle$inf\}<=3, and ||X\textasciicircum{}2-I|$|\_\{inf-\rangle$inf\}<=4 eta are exactly the validated sibling node 1.1.1. Under eta<=1/8, substeps 1.1.2.1--1.1.2.4 prove convergence of B=(X\textasciicircum{}2)\textasciicircum{}(-1/2), B1=1, ||B-I||<=8 eta, S\textasciicircum{}2=I, S1=1, and ||S-X||<=24 eta, so the amended parent statement follows. \steptag{validated} \\[4pt]
\step{1.1.3} With P=theta(X)=(I+S)/2, the identities S\textasciicircum{}2=I and S1=1 give P\textasciicircum{}2=P and P1=1, so P is an exact signed idempotent by def-signed-idempotent; moreover P-Q=(S-X)/2, so ||P-Q|$|\_\{inf-\rangle$inf\} <= C eta, and because Q is entrywise nonnegative row-stochastic by def-stochastic, each negative entry of P is bounded by the corresponding row error, giving delta(P) <= ||P-Q|$|\_\{inf-\rangle$inf\} <= C eta. \steptag{validated} \\[6pt]
\step{1.1.3.1} Dependency import and exact-idempotence algebra: using the prerequisite conclusions of 1.1.1 and 1.1.2, we have X=2Q-I (so Q=(I+X)/2), S\textasciicircum{}2=I, S1=1, and ||S-X|$|\_\{inf-\rangle$inf\}<=24 eta. For P=(I+S)/2, P1=(1+1)/2=1 and P\textasciicircum{}2=(I+2S+S\textasciicircum{}2)/4=(I+S)/2=P, hence P is an exact signed idempotent by def-signed-idempotent. \steptag{archived} \\[4pt]
\step{1.1.3.2} Distance estimate: because P=(I+S)/2 and Q=(I+X)/2, P-Q=(S-X)/2. Therefore ||P-Q|$|\_\{inf-\rangle$inf\}<=(1/2)||S-X|$|\_\{inf-\rangle$inf\}<=12 eta, after the prerequisite bound from 1.1.2 is validated. \steptag{archived} \\[4pt]
\step{1.1.3.3} Negative-mass estimate: by def-stochastic, row-stochastic Q is entrywise nonnegative. For each row i and coordinate j with P\_ij<0, (P\_ij)\_-=-P\_ij<=|P\_ij-Q\_ij| since Q\_ij>=0; summing over j gives the negative mass of row i at most ||p\_i-q\_i||\_1. Taking the maximum over rows gives delta(P)<=||P-Q|$|\_\{inf-\rangle$inf\}<=12 eta. \steptag{archived} \\[4pt]
\step{1.1.3.4} Dependency import and exact-idempotence algebra: using the prerequisite conclusions of 1.1.1 and 1.1.2, we have X=2Q-I (so Q=(I+X)/2), S\textasciicircum{}2=I, S1=1, and ||S-X|$|\_\{inf-\rangle$inf\}<=24 eta. For P=(I+S)/2, P1=(1+1)/2=1 and P\textasciicircum{}2=(I+2S+S\textasciicircum{}2)/4=(I+S)/2=P, hence P is an exact signed idempotent by def-signed-idempotent. \steptag{archived} \\[4pt]
\step{1.1.3.5} Distance estimate: because P=(I+S)/2 and Q=(I+X)/2, P-Q=(S-X)/2. Therefore ||P-Q|$|\_\{inf-\rangle$inf\}<=(1/2)||S-X|$|\_\{inf-\rangle$inf\}<=12 eta, after the prerequisite bound from 1.1.2 is validated. \steptag{archived} \\[4pt]
\step{1.1.3.6} Negative-mass estimate: by def-stochastic, row-stochastic Q is entrywise nonnegative. For each row i and coordinate j with P\_ij<0, (P\_ij)\_-=-P\_ij<=|P\_ij-Q\_ij| since Q\_ij>=0; summing over j gives the negative mass of row i at most ||p\_i-q\_i||\_1. Taking the maximum over rows gives delta(P)<=||P-Q|$|\_\{inf-\rangle$inf\}<=12 eta. \steptag{archived} \\[4pt]
\step{1.1.3.7} Using the prerequisite conclusions of 1.1.1 and 1.1.2: X=2Q-I, so Q=(I+X)/2; Q is row-stochastic and hence Q>=0 entrywise by def-stochastic; S\textasciicircum{}2=I, S1=1, and ||S-X|$|\_\{inf-\rangle$inf\}<=24 eta. With P=theta(X)=(I+S)/2, P1=(1+1)/2=1 and P\textasciicircum{}2=(I+2S+S\textasciicircum{}2)/4=(I+S)/2=P, so P is an exact signed idempotent by def-signed-idempotent. Also P-Q=(S-X)/2, hence ||P-Q|$|\_\{inf-\rangle$inf\}<=12 eta. By the permitted and registered definition def-negative-mass, delta(P)=max\_i sum\_j max\{-P\_ij,0\}. For every i,j, Q\_ij>=0 implies max\{-P\_ij,0\}<=|P\_ij-Q\_ij|: if P\_ij<0 then |P\_ij-Q\_ij|=Q\_ij-P\_ij>=-P\_ij, while if P\_ij>=0 the left side is 0. Summing over j and taking the maximum over i gives delta(P)<=max\_i sum\_j |P\_ij-Q\_ij|=||P-Q|$|\_\{inf-\rangle$inf\}<=12 eta. Thus the node proves exact signed idempotence, the distance bound, and the negative-mass bound with universal constant 12, conditional on validated prerequisites 1.1.1 and 1.1.2. \steptag{validated} \\[6pt]
\step{1.1.3.7.1} Scope/gap certificate for the negative-mass line: this node does not prove delta(P)<=||P-Q|$|\_\{inf-\rangle$inf\} from the current permitted context. In this AF node there are no local dependencies, definitions, or externals, and the lemma workspace registers only def-stochastic and def-signed-idempotent. These inputs give row-stochastic nonnegativity for Q and exact signed idempotence/row-geometry for P, but they do not define delta(P) as max\_i sum\_j max\{-P\_ij,0\}. Therefore the rowwise negative-part argument is unavailable in this scope. The mathematically correct repair is to register the missing negative-mass definition (or weaken the parent statement); until then the parent delta-bound remains unproved. \steptag{validated} \\[6pt]
\step{1.1.3.7.1.1} Allowed-scope audit: this node has no declared dependencies, local definitions, or externals, and the workspace-level af defs are only def-stochastic and def-signed-idempotent. Thus any use of a separate formula for negative mass must come from one of those two allowed definitions or it is out of scope. \steptag{validated} \\[4pt]
\step{1.1.3.7.1.2} The two allowed definitions do not supply the needed formula. def-stochastic gives row-stochastic matrices as entrywise nonnegative maps fixing 1, and def-signed-idempotent gives P1=1 and P\textasciicircum{}2=P together with row-geometry mentioning delta(P). Neither states delta(P)=max\_i sum\_j max\{-P\_ij,0\}. Therefore the expression max\_i sum\_j max\{-P\_ij,0\} cannot be substituted for delta(P) inside this node. \steptag{validated} \\[4pt]
\step{1.1.3.7.1.3} Consequent repair: under the current permitted context the negative-mass estimate is not proved. The parent delta-bound can be discharged only after the missing negative-mass definition is registered as an allowed input, or else the parent statement must be weakened to omit that bound. This child is a gap certificate, not a proof of the delta estimate. \steptag{validated} \\[4pt]
\step{1.1.3.7.2} Scope repair and negative-mass bridge: argument/lemmas/lem-classical-equiv.md permits def-negative-mass, and def-negative-mass is now registered in this AF workspace. Hence delta(P)=max\_i sum\_j max\{-P\_ij,0\} is an allowed definition here, so the old missing-definition caveat is superseded. From 1.1.1, Q is row-stochastic, hence Q\_ij>=0 by def-stochastic. From 1.1.2 and P=(I+S)/2, Q=(I+X)/2, we have P-Q=(S-X)/2 and ||P-Q|$|\_\{inf-\rangle$inf\}<=12 eta. For each i,j, if P\_ij<0 then max\{-P\_ij,0\}=-P\_ij<=Q\_ij-P\_ij=|P\_ij-Q\_ij| because Q\_ij>=0; if P\_ij>=0 the same inequality is trivial. Summing over j and taking max\_i gives delta(P)<=max\_i sum\_j |P\_ij-Q\_ij|=||P-Q|$|\_\{inf-\rangle$inf\}<=12 eta. \steptag{validated} \\[4pt]
\step{1.2} Converse direction: if P is an exact signed idempotent with negative mass delta and Q is obtained by row-normalising the positive parts p\_i\textasciicircum{}+, then Q is row-stochastic, ||P-Q|$|\_\{inf-\rangle$inf\} <= 2 delta, and ||Q\textasciicircum{}2-Q|$|\_\{inf-\rangle$inf\} <= 6 delta+4 delta\textasciicircum{}2. \steptag{validated} \\[6pt]
\step{1.2.1} For row p\_i of P, write p\_i=p\_i\textasciicircum{}+-p\_i\textasciicircum{}- with disjoint nonnegative positive and negative parts, let a\_i=sum\_k (P\_\{ik\})\_- <= delta, and set q\_i=p\_i\textasciicircum{}+/(1+a\_i). Since sum\_k p\_\{ik\}=1 by def-signed-idempotent, sum\_k p\_i\textasciicircum{}+=1+a\_i, so each q\_i is nonnegative of total mass 1 and Q is row-stochastic by def-stochastic. \steptag{validated} \\[4pt]
\step{1.2.2} For e\_i=q\_i-p\_i, one has e\_i=p\_i\textasciicircum{}--(a\_i/(1+a\_i))p\_i\textasciicircum{}+; the two displayed summands have disjoint supports and l1-masses a\_i and a\_i, so ||e\_i||\_1=2a\_i <= 2delta. Therefore D=Q-P satisfies ||D|$|\_\{inf-\rangle$inf\}=max\_i ||e\_i||\_1 <= 2delta, i.e. ||P-Q|$|\_\{inf-\rangle$inf\} <= 2delta. \steptag{validated} \\[4pt]
\step{1.2.3} The row-geometry clause of def-signed-idempotent gives ||P|$|\_\{inf-\rangle$inf\} <= 1+2delta, hence ||P-I|$|\_\{inf-\rangle$inf\} <= 2+2delta. For each row, q\_iQ-q\_i=e\_i(P-I)+q\_iD because P\textasciicircum{}2=P; with q\_i nonnegative of total mass 1, ||q\_iD||\_1 <= ||D|$|\_\{inf-\rangle$inf\} <= 2delta and ||e\_i(P-I)||\_1 <= (2delta)(2+2delta), so ||Q\textasciicircum{}2-Q|$|\_\{inf-\rangle$inf\} <= 6delta+4delta\textasciicircum{}2. \steptag{validated} \\[4pt]
\end{proofbox}

\end{document}
